\ficha{Meissner (2005), A New World Order}{meissner2005}

\begin{identificacao}
MEISSNER, Christopher M. A new world order: explaining the emergence of the
classical gold standard. \emph{Journal of International Economics}, v. 66,
n. 2, p. 385-406, 2005.

Artigo de história econômica quantitativa, em inglês, lido integralmente na
versão de working paper: NBER Working Paper 9233, outubro de 2002, 41 páginas.
As páginas citadas abaixo são as do working paper e não coincidem com a
paginação do artigo publicado.
\end{identificacao}

\bloco{Problema e tese}
Por que cada país aderiu ao padrão-ouro no ano em que aderiu? A difusão foi
gradual, e não simultânea: em 1870 só Grã-Bretanha, Austrália, Canadá e
Portugal tinham padrão-ouro, e em 1910 quase todos tinham (p. 2). A resposta é
que a adesão foi antecipada por externalidades de rede operando pelo comércio,
ou seja, quanto maior o comércio com o bloco-ouro em relação ao PIB, mais cedo
a adesão, e retardada por fragilidade institucional, medida como emissão de
papel alta frente às reservas de ouro, déficit público e dívida. O incentivo de
baratear a captação externa também aparece nos dados. A escolha do padrão não é
tratada como convergência para o sistema mais eficiente, e sim como mudança de
equilíbrio disparada por eventos históricos e depois conduzida por fundamentos
observáveis (p. 9).

\bloco{Argumento}
\begin{enumerate}[leftmargin=*,nosep]
\item Regimes monetários são bens de rede. Compartilhar padrão metálico
economiza custo de transação no comércio, de proteção cambial a comissão de
corretagem na troca de metais, e prata e bimetalismo ofereciam a mesma
propriedade desde que houvesse coordenação (p. 10 e p. 11). Não havia, portanto,
superioridade intrínseca do ouro.
\item O que quebrou o equilíbrio anterior foi a reforma alemã de 1872, que
desencadeou o efeito dominó escandinavo e a limitação da cunhagem de prata na
França (p. 8). O resultado global ser ouro e não prata é creditado à aderência
britânica anterior e ao peso comercial e financeiro da Grã-Bretanha (p. 33).
\item Capacidade institucional condiciona a data. Nenhum país de curso forçado
da amostra aderiu antes de 1878, e a mediana de adesão dos países de papel é
1896, contra 1875 para os já conversíveis em 1870 (p. 14 e p. 15). Finanças
públicas frouxas e bancos mal regulados inviabilizavam a promessa de conversão.
\item O acesso a capital entra como hipótese testável, tomada de Bordo e
Rockoff (1996): se o ouro é selo de boa conduta, país com \emph{spread} alto
sobre o consol britânico tem mais incentivo a aderir depressa (p. 16). Os
exemplos anedóticos citados são Rússia, Japão e China, não a América Latina
(p. 17).
\item Nos resultados, uma elevação de um desvio-padrão na razão comércio com
países-ouro sobre PIB antecipa a adesão em quase 92 por cento do tempo de
espera, e
\chave{This drops the median expected time to adoption from 1879 to about
1871}{p. 24}
Razão de reservas de ouro sobre notas em circulação, déficit sobre produto,
produto per capita e \emph{spread} de títulos têm os sinais previstos e são
significativos (Tabela 3, p. 25).
\item As hipóteses rejeitadas importam tanto quanto as aceitas. Volatilidade
cambial, interesses agrários e devedores rurais, estoque de prata e efeito de
adesão por ideologia não sobrevivem ao teste entre países (p. 27 e p. 28).
\chave{The gold standard appears to be more about transaction costs than
risk-aversion}{p. 34}
\end{enumerate}

\bloco{Conceitos e definições operacionais}
Adotar o padrão-ouro é fixar em lei o preço da moeda nacional em ouro, e só em
ouro, com cunhagem livre e conversibilidade, eventualmente com razão mínima de
reservas para as instituições emissoras (p. 7). A datação exige adesão de
direito e de fato ao mesmo tempo, o que exclui casos de conversibilidade
discricionária como a Áustria-Hungria (p. 9). Externalidade de rede é a queda
de custo de transação proporcional ao tamanho do bloco que usa o mesmo padrão,
com a implicação estratégica de que o padrão do parceiro comercial define a
preferência do país (p. 9 e p. 10). Capacidade institucional é operacionalizada
por três variáveis, razão de reservas de ouro sobre notas emitidas, déficit
sobre PIB e dívida sobre PIB, mais uma \emph{dummy} de regime de papel como
proxy grosseira de responsabilidade fiscal (p. 19 e p. 20). Credibilidade não
recebe definição própria e entra apenas pelo \emph{spread} de títulos.

\bloco{Evidência e método}
Painel desbalanceado de 19 países entre 1870 e 1913, com Brasil, Chile e
Argentina como os latino-americanos (p. 19 e Tabela 6, p. 41). A técnica é
análise de duração, com modelo exponencial de risco na especificação base e
leitura em tempo de falha acelerado, isto é, o coeficiente mede variação
percentual do tempo até a adesão (p. 22 e p. 23). Robustez por Weibull, por
modelo de probabilidade linear com efeitos fixos de país e por variáveis
instrumentais de gravidade, população, distância média aos países-ouro e
fronteira com país-ouro (p. 31). A datação da adesão vem de fontes nacionais
listadas no apêndice, com Calógeras (1960) para o Brasil (p. 40). Não há
contrafactual formal.

\bloco{Agentes e vertentes}
Filiação a Flandreau, que introduziu externalidades de rede na escolha de
regime, e a Eichengreen e Flandreau para a geografia do padrão-ouro e para o
requisito de reservas. Bordo e Rockoff fornecem a hipótese do selo e os dados
de \emph{spread}. Gallarotti responde pela hipótese ideológica, de Cecco pela
disputa entre credores e devedores, Frieden pela leitura da política monetária
americana como interesse exportador. Einaudi e Russell sustentam a narrativa da
Conferência Monetária Internacional de 1867, com Parieu e Lord Liverpool como
agentes nomeados. Sobre o Brasil, a única autoridade citada é Fritsch e Franco
(2000).

\bloco{Críticas e limites}
O resumo afirma que dívida pública alta atrasou a transição, mas o corpo do
texto mostra o coeficiente de dívida sobre produto pequeno, negativo e sem
significância, e a associação só reaparece quando reservas, déficit e produto
per capita saem da regressão (p. 25). A razão de reservas de ouro sobre notas é
reconhecida como possivelmente endógena à intenção de aderir (p. 24). O
\emph{spread} perde significância na especificação longa (p. 27), o que
enfraquece o canal do custo de captação justamente na versão mais completa. E o
Brasil sai da amostra em 1906, ano em que o modelo o dá por aderente: o texto
oferece razões para o Brasil ter chegado ao ouro em 1906 e nada diz sobre o que
se fez depois disso, nem sobre o desenho do arranjo adotado.

\begin{dialogo}
\item Procurar nos editoriais de 1906 o argumento comercial de rede, isto é,
estabilizar o câmbio para reduzir o custo de transação com as praças que já
estavam em ouro. Testa se a defesa da Caixa de Conversão se fez em vocabulário
de comércio exterior e de frete, e não só de preço do café, que é a hipótese
central de Meissner.
\item Procurar como os jornais tratam o efeito esperado sobre o crédito externo
e sobre o juro dos empréstimos. Testa a hipótese do selo de boa conduta, que
Meissner sustenta com Rússia, Japão e China (p. 17) e que ele não documenta
para nenhum país latino-americano.
\item Procurar as menções ao estado do Tesouro e à emissão de papel como
condição prévia da conversibilidade, no vocabulário de época sobre déficit,
lastro e papel-moeda inconversível. Testa a hipótese de capacidade
institucional, que no argumento é o que separa 1906 de qualquer data anterior.
\item Procurar se o argumento da valorização e do interesse do exportador
aparece contra a conversibilidade ou a favor dela. Meissner rejeita, no plano
entre países, que interesses agrários tenham resistido ao ouro (p. 27), e o
caso brasileiro, em que o café pede a Caixa, é o contraexemplo que a regressão
não capta.
\end{dialogo}

\usonocapitulo{Parte I. Serve à afirmação de que a adesão periférica ao
padrão-ouro tem explicação sistemática e não idiossincrática, e que os
determinantes mensuráveis são o comércio com o bloco-ouro, a solidez fiscal e o
custo de captação. É a referência para tratar 1906 como ponto de chegada de um
processo internacional, e não como decisão isolada de política brasileira. O
texto não trata da Caixa de Conversão nem nomeia qualquer agente brasileiro. O
Brasil aparece na Tabela 1 com adesão em 1906 (p. 7), em uma frase sobre bancos
mal regulados e o encilhamento,
\chave{In the 1880s Brazil and Argentina both had weakly regulated banking
systems combined with financially strapped governments meddling in the business
of leading banks}{p. 15}
e em nota de rodapé que registra os papelistas como adversários da
conversibilidade, ao lado dos papeleros chilenos e do Greenback Party (p. 5).
Essa nota é útil na Parte III como registro de que a clivagem imperial
brasileira circulava na literatura internacional, mas é atribuição de segunda
mão, apoiada só em Fritsch e Franco.}
